\section{Exact shift singularities of the entire harmonic difference}
\label{hsh:section}

The Newton expansion in the mixed-spectral report \cite{RepoMixed}
proves holomorphy of its depth-$r$ entire difference on $\Re a>-r$,
and asks for the actual singularity at $a=-r$.  We determine every
negative-integral shift germ, the exact Taylor radius, and the finite
regular term at the first singularity.  The proof uses a finite shift
identity and a cancellation at $a=1$.  Finite transport is a classical
special-function method.  The independent-orders report
\cite{RepoIndependent} supplies a broader entire ordered interpolation
of which this remainder is an exact specialization, identified below.
The contribution here is the explicit singular amplitudes and their
boundary, Taylor-radius, and functional-rank consequences; we do not
assert priority for the underlying transport method.

\subsection{Definitions and the finite shift identity}

For $\Re a>0$, define, initially for $\Re s>1+\Re u$,
\begin{align}
 F_a(s,u)&=\sum_{n\geq0}\frac{(a+u)_n}{(a)_n(n+a)^s},
 &K_a(u)&=\frac{\Gamma(a)\Gamma(1+u)}{\Gamma(a+u)},\label{hsh:definitions}\\
 D(s,u;a)&=F_a(s,u)-K_a(u)F_1(s,u),
 &D_r(s;a)&=[u^r]D(s,u;a).\label{hsh:difference}
\end{align}
Here $(x)_n=x(x+1)\cdots(x+n-1)$ is a rising factorial.  Every
statement in $u$ can be read on a fixed disk $|u|<\rho<1$.
The sources \cite{RepoExact,RepoMixed} prove that $D$ is entire in
$s$, and that its coefficients are jointly holomorphic for
$s\in\mathbb C$, $\Re a>-r$.  In particular,
\begin{equation}\label{hsh:basezero}
 D(s,u;1)=0,\qquad D_0(s;a)=\zeta(s,a)-\zeta(s).
\end{equation}
The pole of the last difference at $s=1$ is removable.

Put
\begin{equation}\label{hsh:P}
 P_n(a,u)=\frac{(a+u)_n}{(a)_n}
       =\sum_{j=0}^n e_j(n;a)u^j,
 \qquad
 e_j(n;a)=\sum_{0\leq\nu_1<\cdots<\nu_j<n}
                  \prod_{i=1}^j\frac1{a+\nu_i}.
\end{equation}
The empty product is one and $e_j(n;a)=0$ for $j>n$.

\paragraph{Exact identification with the ordered endpoint interpolation.}
Let $\mathcal H_{b,a}(w)$ denote the relative character in
\cite{RepoIndependent}, normalized by the deconcatenation identity
$Z_a(w)=\sum_{w=vw'}Z_b(v)\mathcal H_{b,a}(w')$ for the
strict ordered Hurwitz sums.  Its Gamma formula for words of ones is
$\mathcal H_{1,a}(\{1\}^j)=[u^j]K_a(u)$.  Applying the defining
recursion to $(s,\{1\}^r)$ therefore gives the exact identification
\begin{equation}\label{hsh:relativecharacter}
          D_r(s;a)=\mathcal H_{1,a}(s,\{1\}^r).
\end{equation}
Indeed $Z_a(s,\{1\}^r)=\sum_{n\geq0}e_r(n;a)(n+a)^{-s}$,
and the other deconcatenation terms are exactly the coefficients
of $K_a(u)F_1(s,u)$ in \eqref{hsh:difference}.  This first proves
the identity in a convergence domain, then everywhere by the entire
spectral continuations.  Proposition~\ref{hsh:transport} is accordingly
also a specialization of that source's endpoint transport.

\begin{proposition}[Exact finite transport]\label{hsh:transport}
For every integer $L\geq1$,
\begin{equation}\label{hsh:finite}
 D(s,u;a)=\sum_{n=0}^{L-1}P_n(a,u)(a+n)^{-s}
                       +P_L(a,u)D(s,u;a+L).
\end{equation}
Consequently
\begin{align}
 D(s,u;a)&=a^{-s}+\frac{a+u}{a}D(s,u;a+1),
                                                     \label{hsh:one}\\
 D_r(s;a)&=\mathbf1_{r=0}a^{-s}+D_r(s;a+1)
                      +\mathbf1_{r\geq1}\frac{D_{r-1}(s;a+1)}a,
                                                     \label{hsh:depthshift}\\
 D_r(s;a)&=\sum_{n=r}^{L-1}e_r(n;a)(a+n)^{-s}
          +\sum_{j=0}^{\min(r,L)}e_j(L;a)D_{r-j}(s;a+L).
                                                     \label{hsh:depthfinite}
\end{align}
The finite sum in the last formula is empty when $L\leq r$.
These identities hold wherever the powers and continuations on both
sides use the same branches.
\end{proposition}

\begin{proof}
Separating the $n=0$ summand gives
$F_a=a^{-s}+(a+u)F_{a+1}/a$.  The Gamma functional equation gives
$K_a=(a+u)K_{a+1}/a$.  Subtraction proves
\eqref{hsh:one} in the initial common half-plane.  Both sides are
entire in $s$, so that identity continues to every $s$.  Iterating it
gives \eqref{hsh:finite}; coefficient extraction gives the other
formulas.  All operations are finite after the initial identity, so
they create no interchange of a divergent sum and a derivative.
\end{proof}

\subsection{Every negative-integral shift germ}

For $N\geq0$, put $a=-N+z$ and define
\begin{equation}\label{hsh:C}
 C_N(z,u)=P_N(-N+z,u)
          =\prod_{j=1}^{N}\left(1+\frac{u}{z-j}\right),
 \qquad c_{N,r}(z)=[u^r]C_N(z,u).
\end{equation}
Thus $C_0=1$ and $c_{N,r}=0$ if $r>N$.  Near $z=0$ choose
a branch of $\log z$.  For the other powers in
\eqref{hsh:finite}, choose logarithms of $z-h$, $1\leq h\leq N$,
holomorphic on a small disk around zero.  Their values at zero may be
$\log h+\mathrm i\pi$ or $\log h-\mathrm i\pi$ for the two lateral
continuations.

\begin{theorem}[Exact local decomposition]\label{hsh:germ}
On these specified branches,
\begin{equation}\label{hsh:local}
 D(s,u;-N+z)=C_N(z,u)z^{-s}+H_N(s,u;z),
\end{equation}
where $H_N$ is jointly holomorphic in $s$, $u$, and $z$ near
every point $(s_0,0,0)$.  An explicit formula is
\begin{align}
 H_N(s,u;z)={}&
 \sum_{n=0}^{N-1}P_n(-N+z,u)(z-N+n)^{-s}\notag\\
 &+\left(1+\frac uz\right)C_N(z,u)D(s,u;1+z).
                                                     \label{hsh:H}
\end{align}
The apparent pole in the second line is removable.  For every
$r,k\geq0$,
\begin{equation}\label{hsh:localjets}
 \partial_s^kD_r(s;-N+z)
   =c_{N,r}(z)z^{-s}(-\log z)^k
                         +\partial_s^kH_{N,r}(s;z),
 \qquad H_{N,r}=[u^r]H_N.
\end{equation}
For $r\leq N$ the leading coefficient is nonzero:
\begin{equation}\label{hsh:leading}
 c_{N,r}(0)=(-1)^r e_r\left(1,\frac12,\ldots,\frac1N\right).
\end{equation}
For $r>N$ the entire singular amplitude vanishes identically.
\end{theorem}

\begin{proof}
Use \eqref{hsh:finite} with $L=N+1$.  Its $n=N$ term is exactly
$C_N(z,u)z^{-s}$, and the other terms give \eqref{hsh:H} because
$P_{N+1}(-N+z,u)=(1+u/z)C_N(z,u)$.
Every factor of $P_n$, $n\leq N$, is holomorphic in $z$ near
zero.  The chosen logarithms make all powers with $n<N$
holomorphic there as well.  By \eqref{hsh:basezero},
$D(s,u;1+z)=z\,\widetilde D(s,u;z)$ with $\widetilde D$
jointly holomorphic.  This cancels the sole $1/z$ factor in the
last line.  This also proves local normality on compact spectral
sets, so all the stated derivatives can be taken.  Finally,
$C_N(0,u)=\prod_{j=1}^{N}(1-u/j)$, which proves
\eqref{hsh:leading}; its elementary symmetric coefficients are
strictly positive before the displayed sign.
\end{proof}

\begin{corollary}[Principal continuation]\label{hsh:principal}
For each $r\geq0$, $D_r(s;a)$ has a jointly holomorphic principal
continuation to
\begin{equation}\label{hsh:slit}
       (s,a)\in\mathbb C\times
                         \bigl(\mathbb C\setminus(-\infty,-r]\bigr).
\end{equation}
It is holomorphic across $a=0,-1,\ldots,-(r-1)$.  On the
specified lateral branches, every $a=-N$ with $N\geq r$ has
the germ \eqref{hsh:localjets}.
\end{corollary}

\begin{proof}
On any compact set in the slit plane, choose $L$ so large that
$\Re(a+L)>0$.  In \eqref{hsh:depthfinite}, the only powers that
remain have $n\geq r$, hence their arguments avoid the principal
logarithm cut.  The tail is defined in the original half-plane.
The rational coefficients have only possible poles at negative
integers.  Those in the slit domain have $N<r$, and
Theorem~\ref{hsh:germ} proves their removability.  The formulas
agree on overlaps by finite transport and coincide with the
original function on a common open set.  They therefore give
the claimed jointly holomorphic continuation.
\end{proof}

\paragraph{A distinction between principal continuation and all sheets.}
The early removable points cannot simply be deleted from the
puncture set for every sheet.  A positive circuit around $a=-N$
changes the chosen germ by
\begin{equation}\label{hsh:monodromy}
 \Delta_N D_r(s;a)=
 (\exp(-2\pi\mathrm i s)-1)e_r(N;a)(a+N)^{-s}.
\end{equation}
This follows by changing only $\log(a+N)$ in a finite transport
formula with $L>N$.  The rational coefficient $e_r(N;a)$ may have
poles at earlier shifts.  For instance, a circuit around $-1$
changes $D_1(s;a)$ by
\[
       (\exp(-2\pi\mathrm i s)-1)\frac{(a+1)^{-s}}a.
\]
Although the principal $D_1$ is regular at $a=0$, this continued
branch has a pole there for generic $s$.  At $s=-m$, differentiating
once in $s$ gives the change $-2\pi\mathrm i(a+1)^m/a$.
Thus the same issue occurs for spectral derivatives at the polynomial
orders.  Formula \eqref{hsh:slit} is a statement about the principal
branch.  Arbitrary multivalued continuation can safely be made along
paths avoiding all of $0,-1,-2,\ldots$; it need not preserve the
removability of every preimage of an earlier shift.  This is a
continuation qualification, not a counterexample to the source's
half-plane theorem.

\subsection{The first boundary and Arakawa--Kaneko values}

For $j\geq1$, write
\begin{equation}\label{hsh:Xi}
 \Xi_j(s)=\frac1{\Gamma(s)}\int_0^\infty
   \frac{t^{s-1}}{e^t-1}\operatorname{Li}_j(1-e^{-t})\,dt,
 \qquad \Re s>0.
\end{equation}
These are the classical Arakawa--Kaneko functions, continued entirely
in $s$ by Taylor subtraction at zero
\cite{ArakawaKaneko,KanekoTsumura}.  In the descending convention,
the multiple version $\Xi(j,\{1\}^{p-1};s)$ is obtained by replacing
the polylogarithm in \eqref{hsh:Xi} by
$\operatorname{Li}_{j,1,\ldots,1}$.
The mixed-spectral report proves the exact shift-jet identity
\begin{equation}\label{hsh:AKjet}
 \left.\partial_a^pD_q(s;a)\right|_{a=1}
      =(-1)^{q+p}p!\,\Xi(q+1,\{1\}^{p-1};s),\qquad p\geq1.
\end{equation}
In particular $\Xi_1(s)=s\zeta(s+1)$, with its removable value
at $s=0$ understood.  This follows directly from
$\operatorname{Li}_1(1-e^{-t})=t$.

For a fixed positive depth $r$, set
\begin{align}
 E_h^{(r)}&=e_h\left(1,\frac12,\ldots,\frac1r\right),
 \qquad\prod_{j=1}^r\left(1-\frac zj\right)^{-1}
       =\sum_{n\geq0}h_n^{(r)}z^n,\label{hsh:eh}\\
 c_r(z)&=c_{r,r}(z)=\frac1{\prod_{j=1}^r(z-j)}
       =\frac{(-1)^r}{r!}\prod_{j=1}^r\left(1-\frac zj\right)^{-1},\notag\\
 B_r(s)&=(-1)^r\sum_{h=0}^{r-1}E_h^{(r)}\Xi_{r-h}(s).
                                                       \label{hsh:B}
\end{align}
The $h_n^{(r)}$ are complete homogeneous symmetric polynomials.

\begin{theorem}[Exact first-boundary normalization]\label{hsh:boundary}
For $r\geq1$,
\begin{equation}\label{hsh:boundaryidentity}
 \boxed{\displaystyle
 \lim_{z\to0}\left\{D_r(s;-r+z)-c_r(z)z^{-s}\right\}=B_r(s).}
\end{equation}
The expression inside braces extends holomorphically through $z=0$.
Convergence is locally uniform in $s$, and arbitrary fixed spectral
derivatives can be applied to the identity.  In particular,
\begin{align}
 B_1(s)&=-s\zeta(s+1),\notag\\
 B_2(s)&=\Xi_2(s)+\tfrac32s\zeta(s+1),\notag\\
 B_3(s)&=-\Xi_3(s)-\tfrac{11}{6}\Xi_2(s)-s\zeta(s+1).
                                                       \label{hsh:examples}
\end{align}
Every Taylor coefficient of the holomorphic remainder in
\eqref{hsh:boundaryidentity} is a finite combination of the multiple
functions in \eqref{hsh:AKjet}, with rational coefficients.
\end{theorem}

\begin{proof}
Take $N=r$ in \eqref{hsh:H} and extract $u^r$.  The finite terms
with $n<r$ vanish because $P_n$ has degree $n$.  Thus the remainder
is exactly
\begin{equation}\label{hsh:firstH}
 H_{r,r}(s;z)=[u^r]\left(1+\frac uz\right)
                          C_r(z,u)D(s,u;1+z).
\end{equation}
At $z=0$, the term without $u/z$ vanishes.  The remaining term is
$[u^r]uC_r(0,u)\partial_aD(s,u;1)$.
Now $C_r(0,u)=\sum_{h=0}^r(-1)^hE_h^{(r)}u^h$, while
\eqref{hsh:AKjet} at $p=1$ gives
$\partial_aD(s,u;1)=\sum_{q\geq0}(-1)^{q+1}\Xi_{q+1}(s)u^q$.
Their finite coefficient product is precisely $B_r$.
The holomorphy assertion is already contained in
Theorem~\ref{hsh:germ}.  For higher Taylor coefficients, expand
$C_r(z,u)$ rationally at $z=0$ and use
\[
 D(s,u;1+z)=\sum_{p\geq1}\sum_{q\geq0}
       (-1)^{p+q}\Xi(q+1,\{1\}^{p-1};s)u^qz^p.
\]
At a fixed power of $z$ in \eqref{hsh:firstH}, only finitely many
$p,q$ occur.  This also proves the final assertion with its
normalization specified.
\end{proof}

The same calculation gives every regular boundary constant.  With
$E_h^{(N)}=e_h(1,1/2,\ldots,1/N)$ and $E_0^{(0)}=1$,
\begin{equation}\label{hsh:allregular}
 \begin{split}
 H_{N,r}(s;0)={}&\sum_{n=0}^{N-1}e_r(n;-N)(n-N)^{-s}\\
 &+(-1)^r\sum_{h=0}^{\min(N,r-1)}E_h^{(N)}\Xi_{r-h}(s).
 \end{split}
\end{equation}
The second sum is empty for $r=0$.  This follows by evaluating the
first line of \eqref{hsh:H} at zero and extracting the coefficient
of $u^r$ from $uC_N(0,u)\partial_aD(s,u;1)$ in its second line.
The finite powers retain their specified lateral logarithms.  If
$r>N$, the first sum is empty and the singular amplitude vanishes;
the formula recovers the source's previously proved negative-shift
values.  If $r=N\geq1$, it is \eqref{hsh:boundaryidentity}.

\paragraph{The regular boundary term and a coordinate finite part.}
The subtraction in \eqref{hsh:boundaryidentity} retains the full
analytic amplitude $c_r(z)$; it is not the same operation as discarding
only the divergent terms of a Laurent expansion.  For $p\geq0$,
take the finite part along the positive real coordinate $z$ to mean
the coefficient of $z^0(\log z)^0$ in the power--logarithm expansion.
Then
\begin{equation}\label{hsh:FP}
 \boxed{\displaystyle
 \FP_{z=0}D_r(p;-r+z)=B_r(p)+\frac{(-1)^r}{r!}h_p^{(r)}.}
\end{equation}
Indeed the constant contributed by $c_r(z)z^{-p}$ is
$[z^p]c_r(z)$.  On the other hand, for every $k\geq1$,
\begin{equation}\label{hsh:FPjets}
 \FP_{z=0}\left.\partial_s^kD_r(s;-r+z)\right|_{s=p}
                       =B_r^{(k)}(p),\qquad p\geq0,
\end{equation}
because its singular amplitude has the factor $(-\log z)^k$.
These formulas concern a fixed coordinate and a defined constant
coefficient, so there is no implicit conversion between two
regularization prescriptions.

For $m\geq1$ and $k\geq0$, there is an ordinary lateral limit,
\begin{equation}\label{hsh:negativeboundary}
 \lim_{z\to0}\left.\partial_s^kD_r(s;-r+z)\right|_{s=-m}
                      =B_r^{(k)}(-m),
\end{equation}
provided $\arg z$ stays in a bounded interval on the chosen branch.
The reason is $z^m(\log z)^k\to0$; no finite part is needed.
For positive $k$, this continuous extension is not
holomorphic, as the next subsection proves.

\paragraph{Explicit Stieltjes and Gamma boundary identities at depth one.}
Use the convention
$\zeta(1+s)=s^{-1}+\sum_{j\geq0}(-1)^j\gamma_js^j/j!$.
Then \cite{DLMFHurwitz}
\begin{equation}\label{hsh:Stieltjes}
 B_1(0)=-1,\qquad
 B_1^{(k)}(0)=(-1)^k k\gamma_{k-1}\quad(k\geq1).
\end{equation}
Consequently, for example,
\begin{align}
 \lim_{z\to0}\{D_1(0;-1+z)-(z-1)^{-1}\}&=-1,\notag\\
 \lim_{z\to0}\{\partial_sD_1(0;-1+z)
                          +(z-1)^{-1}\log z\}&=-\gamma,\notag\\
 \lim_{z\to0}\{\partial_s^2D_1(0;-1+z)
                          -(z-1)^{-1}(\log z)^2\}&=2\gamma_1.
                                                       \label{hsh:stieltjeslimits}
\end{align}
At the negative orders, for $m\geq1$,
\begin{equation}\label{hsh:zetaexplicit}
 B_1^{(k)}(-m)=m\zeta^{(k)}(1-m)
                         -k\zeta^{(k-1)}(1-m),
\end{equation}
where the last term is omitted for $k=0$.  In particular,
\[
 B_1'(-1)=\frac12-\frac12\log(2\pi),\qquad
 B_1'(-2)=\frac1{12}+2\zeta'(-1).
\]
The first evaluation uses the classical Lerch normalization of
$\zeta'(0)$ \cite{DLMFHurwitz}.  These constants occur here as
boundary values of a continued harmonic family, with the complete
singular amplitude specified.

\subsection{Normalized Lerch and polylogarithmic primitives}
\label{hsh:primitives}

The exact singular amplitude also gives a finite antiderivative formula.
This is an application of standard partial fractions and Lerch calculus
\cite{DLMFLerch}, with the normalization needed at resonant spectral
orders made explicit.  Fix a logarithm on a slit punctured unit disk
and a base point $z_0$ in that disk.  For $1\leq j\leq r$, put
\begin{align}
 A_{r,j}&=\frac{(-1)^{r-j}}{(j-1)!(r-j)!},
 &c_r(z)&=\sum_{j=1}^{r}\frac{A_{r,j}}{z-j},
                                                    \label{hsh:primitivePF}\\
 P_j(s;z,z_0)&=-\frac1j\left\{
 z^{1-s}\Phi(z/j,1,1-s)
       -z_0^{1-s}\Phi(z_0/j,1,1-s)\right\}.
                                                    \label{hsh:Lerchprimitive}
\end{align}
Here $\Phi(w,1,v)=\sum_{n\geq0}w^n/(n+v)$ for $|w|<1$,
continued meromorphically in $v$.

\begin{proposition}[Entire normalized spectral primitive]
\label{hsh:primitive}
The difference $P_j$ extends to an entire function of $s$, including
the apparent resonances $s=1,2,\ldots$, and
\begin{equation}\label{hsh:primitivederivative}
 P_j(s;z_0,z_0)=0,\qquad
 \frac{\partial}{\partial z}P_j(s;z,z_0)=\frac{z^{-s}}{z-j}.
\end{equation}
Consequently a normalized primitive of the complete first-boundary germ is
\begin{equation}\label{hsh:fullprimitive}
 \sum_{j=1}^{r}A_{r,j}P_j(s;z,z_0)
                   +\int_{z_0}^{z}H_{r,r}(s;t)\,\mathrm dt.
\end{equation}
The remainder in this formula is holomorphic on the unit disk.
Every fixed spectral derivative can be taken in the formula.
At $s=0$, all the differentiated singular primitives are finite
ordinary-polylogarithm expressions: define
\begin{equation}\label{hsh:polylogprimitive}
 J_{j,k}(z)=-k!\sum_{h=0}^{k}\frac{(-\log z)^h}{h!}
                          \operatorname{Li}_{k+1-h}(z/j).
\end{equation}
Then
\begin{equation}\label{hsh:polylogprimitivejet}
 \left.\partial_s^kP_j(s;z,z_0)\right|_{s=0}
                         =J_{j,k}(z)-J_{j,k}(z_0),
 \qquad J_{j,k}'(z)=\frac{(-\log z)^k}{z-j}.
\end{equation}
\end{proposition}

\begin{proof}
The residues of $c_r$ at its distinct simple poles give
\eqref{hsh:primitivePF}.  Expanding the Lerch functions gives the
normally convergent, already normalized series
\begin{equation}\label{hsh:primitiveentire}
 P_j(s;z,z_0)=
 -\sum_{m\geq1}\frac{z^{m-s}-z_0^{m-s}}{j^m(m-s)}.
\end{equation}
For a term with $s=m$, interpret the quotient as
$\log z-\log z_0$.  Each term is entire in $s$.  On compact
spectral sets and compact subsets of the slit punctured disk, its tail
is bounded by a geometric series times $1/m$.  Hence the sum and all
fixed spectral derivatives are locally normal.  Equivalently, each
unnormalized Lerch primitive in \eqref{hsh:Lerchprimitive} has residue
$j^{-p}$ at $s=p\geq1$, independent of its endpoint; the difference
cancels that residue.  Differentiating \eqref{hsh:primitiveentire}
in $z$ proves \eqref{hsh:primitivederivative}.  Formula
\eqref{hsh:firstH} shows that $H_{r,r}$ is holomorphic for $|z|<1$,
with its sole apparent pole at zero removed.  Thus its integral is
well defined independently of the path in that disk, and
\eqref{hsh:fullprimitive} follows.  Finally, differentiating
$z^{-s}/(m-s)$ at $s=0$ and summing the resulting powers $m^{-q}$
gives \eqref{hsh:polylogprimitive}; termwise differentiation, or the
usual polylogarithm derivative rule, proves
\eqref{hsh:polylogprimitivejet}.
\end{proof}

In particular,
\[
 J_{j,0}(z)=\log(1-z/j),\qquad
 J_{j,1}(z)=-\operatorname{Li}_2(z/j)
                        -\log z\log(1-z/j).
\]
The sign in the second formula is fixed by
$J_{j,1}'(z)=-\log z/(z-j)$.  The base-point subtraction supplies
the additive constant before continuation or spectral differentiation;
separate unnormalized Lerch terms should not be evaluated at their
resonant parameter poles.

\subsection{Exact Taylor radii and the rank jump at integral orders}

\begin{theorem}[Exact shift Taylor radius]\label{hsh:radius}
Fix a real center $a_*>-r$ and a spectral order $s_0$.
The Taylor series in $a-a_*$ of $D_r(s_0;a)$ has radius
\[
 \begin{cases}
 a_*+r,&s_0\notin\{0,-1,-2,\ldots\},\\
 \infty,&s_0\in\{0,-1,-2,\ldots\}.
 \end{cases}
\]
For every $k\geq1$ and every $s_0\in\mathbb C$, the Taylor
series of $\partial_s^kD_r(s_0;a)$ has radius exactly $a_*+r$.
In particular, at the center $a_*=1$, the lower bound $r+1$
in \cite{RepoMixed} is exact outside the polynomial cases, and is
exact for every positive spectral derivative even at those cases.
\end{theorem}

\begin{proof}
The disk $|a-a_*|<a_*+r$ lies in the domain
\eqref{hsh:slit}, so holomorphy gives the lower bound, uniformly
on compact spectral sets.  At $a=-r$, the singular amplitude is
$c_r(z)z^{-s_0}(-\log z)^k$, with $c_r(0)=(-1)^r/r!\ne0$;
for $r=0$ the same statement has $c_0=1$.
When $k=0$ and $s_0$ is not a nonpositive integer, $z^{-s_0}$
has either a pole or nontrivial branching.  It therefore cannot
extend holomorphically to zero.  For $k\geq1$, its logarithmic
factor prevents holomorphic extension for every $s_0$.  More
explicitly, one circuit sends the germ to
$e^{-2\pi\mathrm i s_0}z^{-s_0}(-\log z-2\pi\mathrm i)^k$.
If $s_0$ is nonintegral, the exponential factor is nontrivial;
if $s_0$ is integral, the nonconstant logarithmic polynomial is
not translation invariant.  Positive integral $s_0$ at $k=0$
instead gives a pole.  The holomorphic remainder cannot cancel
any of these obstructions.  Hence the Taylor disk cannot be
larger.  At $s_0=-m$, the finite Stirling formula below proves
that $D_r(-m;a)$ is a polynomial and has infinite radius.
\end{proof}

\begin{corollary}[Functional independence of spectral jets]
\label{hsh:independence}
If $s_0\notin\{0,-1,-2,\ldots\}$, every finite family
\[
 \{\partial_s^kD_r(s_0;a):0\leq r\leq R,\ 0\leq k\leq K\}
\]
is linearly independent over $\mathbb C$, even modulo entire
functions of $a$.  If $s_0=-m$ with $m\geq0$, the same assertion
holds for the subfamily $1\leq k\leq K$.
This is independence of functions of the shift, not arithmetic
independence of any special values.
\end{corollary}

\begin{proof}
Suppose a finite linear combination equals an entire function,
and let $r$ be the smallest depth with a nonzero coefficient.
Every larger-depth term is holomorphic at $a=-r$ on the principal
continuation.  The singular germ of the depth-$r$ terms is
$c_r(z)z^{-s_0}P(-\log z)$ for a nonzero polynomial $P$.
The monodromy argument from Theorem~\ref{hsh:radius} shows that
this cannot be holomorphic, except possibly when $s_0$ is a
nonpositive integer and $P$ is constant.  That exception is
excluded by the stated subfamily.  Thus all coefficients at
this depth vanish, a contradiction.  Induction gives the claim.
\end{proof}

Let $\left\{\begin{smallmatrix}n\\j\end{smallmatrix}\right\}$
denote a Stirling number of the second kind and let
$(a-1)^{\underline j}=(a-1)\cdots(a-j)$.
The known finite formula \cite{RepoExact,RepoMixed} is
\begin{equation}\label{hsh:finiteStirling}
 D_r(-m;a)=\sum_{j=1}^{m+1}
 \left\{\begin{matrix}m+1\\j\end{matrix}\right\}
           (a-1)^{\underline j}\left(-\frac1j\right)^{r+1}.
\end{equation}

\begin{proposition}[Exact finite rank and minimal depth recurrence]
\label{hsh:rank}
For $m\geq0$, the functions $\{D_r(-m;a):r\geq0\}$ span
exactly the $(m+1)$-dimensional space of polynomials of degree
at most $m+1$ vanishing at $a=1$.  The first $m+1$ depths form
a basis.  If $E$ shifts depth, $ED_r=D_{r+1}$, then
\begin{equation}\label{hsh:recurrence}
 \boxed{\displaystyle
       \prod_{j=1}^{m+1}\left(E+\frac1j\right)D_r(-m;a)=0,
       \qquad r\geq0.}
\end{equation}
Its order $m+1$ is minimal among nonzero constant-coefficient
depth recurrences valid as identities in $a$.
For any fixed $k\geq1$, the family
$\{\partial_s^kD_r(-m;a):r\geq0\}$ is infinite dimensional and
satisfies no nonzero constant-coefficient depth recurrence.
The latter statement also holds at nonexceptional $s_0$ with $k=0$.
\end{proposition}

\begin{proof}
The polynomials
$\left\{\begin{smallmatrix}m+1\\j\end{smallmatrix}\right\}
(a-1)^{\underline j}$, $1\leq j\leq m+1$, have distinct degrees,
nonzero leading coefficients, and vanish at one.  They form a basis
of the stated space.  The coefficient matrix for depths $0,\ldots,m$
in \eqref{hsh:finiteStirling} has entries $\lambda_j^{r+1}$,
where $\lambda_j=-1/j$.  Its determinant is a nonzero product
of the $\lambda_j$ and a Vandermonde determinant.  This proves
the rank and basis assertions.  Each component in
\eqref{hsh:finiteStirling} is annihilated by $E-\lambda_j$,
which proves \eqref{hsh:recurrence}.  Conversely, independence
of those polynomial components forces an annihilating polynomial
$Q$ to vanish at every $\lambda_j$.  Its degree is therefore
at least $m+1$.  Finally any recurrence for a positive spectral
derivative would give a finite constant linear relation among
distinct depths, contradicting Corollary~\ref{hsh:independence}.
\end{proof}

For example, at $s=0$ one has
$D_r(0;a)=(-1)^{r+1}(a-1)$ and $(E+1)D_r=0$.
At $s=-1$ the minimal recurrence is
$D_{r+2}+\tfrac32D_{r+1}+\tfrac12D_r=0$.
A single spectral differentiation removes these finite-rank
relations.  This failure is detected by exact logarithmic germs,
without any numerical independence assumption.

\subsection{Audit and reproducibility}

The source's holomorphy domain and Taylor lower bound were correct;
the new statements determine their actual boundary behavior.  The
regular boundary term $B_r$ and the coordinate finite part in
\eqref{hsh:FP} must retain their different definitions.  Likewise,
principal-sheet removability must not be promoted to removability on
all sheets: \eqref{hsh:monodromy} supplies an explicit obstruction.

The accompanying program \src{code/verify_shift_germs.py} checks the
finite recurrence and residue algebra exactly with rational arithmetic,
compares the Newton series with the independently convergent defining
harmonic sums at complex parameters, checks finite transport at negative
shifts, and compares Cauchy-extracted boundary constants with
Arakawa--Kaneko integrals and ordinary zeta derivatives.  The numerical
checks are finite-precision diagnostics with separately recorded
truncations; they are not interval certificates or replacements for
the normal-convergence and continuation arguments above.
The short companion \src{code/verify_shift_primitives.py} checks the
partial-fraction amplitudes and the polylogarithmic primitive derivatives
with exact symbolic arithmetic.
